% =============================================================================
\section{Introduction}
% =============================================================================

\subsection{Motivation}

A photograph is a two-dimensional projection of a three-dimensional world, and
recovering that world from a single view is underconstrained: infinitely many
scenes project to the same image.
Humans nonetheless read a photograph as a place. We infer that the ground
continues past the frame, that the occluded side of a rock exists, that the
flowers in the foreground are individual plants rather than a printed texture.
Building an environment that supports those inferences, one a person can
\emph{jump into}, is the problem this project begins to address.

Photographs are among the most abundant records of place we have: personal
archives, cultural heritage documentation, scientific field imagery, and now,
with generative models, images of places that never existed at all. A system that
converts such an image into an environment a person can walk into changes what
those archives are able to express. It also lowers the cost of 3D content
authoring by several orders of magnitude relative to manual modeling or
multi-view capture, and because the process is automatic, it puts that authoring
within reach of a non-expert.

\subsection{The Challenge}

The difficulty is not one hard sub-problem but the interaction of several, each
of which degrades the others:

\begin{description}[leftmargin=1.2em,style=nextline]
  \item[Extreme ill-posedness.] A single view fixes neither absolute scale nor
  the geometry of anything it does not see. Roughly $5/6$ of a full
  $360^\circ$ panorama generated from a standard photograph is synthesized rather
  than observed, and everything behind every foreground object is unknown. Much
  of the world must be inferred.

  \item[Compounding error.] A pipeline that estimates depth, then geometry, then
  materials, then objects inherits every earlier error, and a discrepancy that is
  barely visible at one stage can be structural by the next.

  \item[Heterogeneous content.] The representation that suits a mountainside
  (a height field) is wrong for a tree (a mesh instance), which is wrong for a
  meadow (a scattered population), which is wrong for the sea (an animated
  surface). These representations must nonetheless coexist in one scene: they
  need to share a coordinate frame and a scale, agree on how they are lit, and
  not seam visibly where they meet.

  \item[Interactivity constraints.] The output must render stereoscopically at
  headset frame rates. That budget constrains every rendering choice, and it
  binds hardest on the animated, simulated elements that make the scene
  interactive in the first place.

\end{description}

\subsection{Goals}

Our goal is a system that takes one image and produces an environment which is
\emph{interactive} (objects animate and interact with their surroundings),
\emph{decomposed} (discrete objects exist as individual entities, not painted
texture), \emph{plausible} beyond the observed region, and \emph{performant} at
VR frame rates.
